%
% Front matter
%

\begin{frontmatter}

\title{
  T2K 6.3933$\times10^{20}$-POT 3-Flavour Muon-Neutrino Disappearance Analysis
}

\vspace*{0.7cm}

\author[RAL]     {Costas Andreopoulos}
\author[RAL,OXF] {Thomas Dealtry}
\author[IFIC]    {Lorena Escudero S\'anchez}
\author[LANCS]   {Nick Grant}
\author[RAL,OXF] {and Alfons Weber}

\address[RAL]{Particle Physics Dept., Rutherford Appleton Laboratory, STFC \\
              Harwell Oxford Campus\\
              Oxfordshire OX11 0QX, UK}

\address[OXF]{Subdepartment of Particle Physics, University of Oxford \\
              Denys Wilkinson Building, Keble Road\\
              Oxfordshire OX1 3RH, UK}

\address[IFIC]{Instituto de Fisica Corpuscular\\
              CSIC and University of Valencia\\
              Valencia, Spain}

\address[LANCS]{Physics Dept., Lancaster University \\
              Lancaster, LA1 4YB, UK}

%
% Abstract
%

\begin{abstract}

We report the results of the
3-flavour $\nu_{\mu}$-disappearance analysis,
performed by the {\em VaLOR} (Valencia-Lancaster-Oxford-Rutherford) group,
on the combined Run 1+2+3+4a dataset.
This analysis is an update of the work presented in
Refs. \cite{OurRun1DisappAnalysis}, \cite{OurRun12DisappAnalysis}, 
\cite{OurRun12DisappAnalysis3Flavour} and \cite{OurRun123DisappAnalysis3Flavour}
with improved methods and an enlarged dataset.
The Run 1+2+3+4a dataset corresponds to an integrated J-PARC neutrino beam
exposure of 6.3933$\times$10$^{20}$ POT.

This analysis predicts 421.76 $\pm$ 23.33 (syst)
1 $\mu$-like ring events in SuperK in the absence of any oscillation,
but only 115 were observed. The observed deficit has a strong energy dependence;
the ratio of observed to expected, under the no-oscillation hypothesis, is
$\sim$27\% $<$ 0.5 GeV, 
$\sim$10\% between 0.5 and 1 GeV and 
$\sim$73\% $>$ 1 GeV.

A $\nu_{\mu}$-disappearance analysis was performed 
in a framework of 3-flavour oscillations including matter effects in constant-density matter.
The observed reconstructed energy spectrum of 1 $\mu$-like ring events was fitted
allowing $\sin^{2}\theta_{23}$  and $|\Delta m^{2}_{32}|$ to float
$\delta_{CP}$ was fixed at 0 and a mass hierarchy was
assumed\footnote{
  Two independent fits were performed. The first assumed the normal mass hierarchy.
  The second assumed the inverted mass hierarchy.
}.
The other \emph{non-23-sector} oscillation parameters,
and all 41 systematic parameters considered in this analysis, were also allowed to float in the fit.

The 3-flavour $\nu_{\mu}$-disappearance fit to the combined Run 1+2+3+4a dataset,
assuming the normal mass hierarchy,
%with $\sin^{2}\theta_{23}$ constrained to the physical region,
yields
$|\Delta m^{2}_{32}|$ = 2.463$\times$10$^{-3}$ eV$^{2}$/c$^{4}$  and
$\sin^{2}\theta_{23}$ = 0.514
as the best-fit values (p-value = 0.88).
The fit assuming the inverted mass hierarchy yields
$|\Delta m^{2}_{32}|$ = 2.530$\times$10$^{-3}$ eV$^{2}$/c$^{4}$  and
$\sin^{2}\theta_{23}$ = 0.511
as the best-fit values (p-value = 0.89).

Confidence regions in $\sin^{2}\theta_{23}$ - $|\Delta m^{2}_{32}|$ space 
are shown using both the constant-$\Delta\chi^2$ method and the method of Feldman and Cousins.
Using the profile $\Delta\chi^2$ distributions obtained from the fit
%and the method of Feldman and Cousins, 
the 90\% (68\%) CL allowed regions for $|\Delta m^{2}_{32}|$ and $\sin^{2}\theta_{23}$
can be summarized as
2.30$\times$10$^{-3}$ (2.35$\times$10$^{-3}$) eV$^{2}$/c$^{4}$ 
$< |\Delta m^{2}_{32}| < 
$ 2.63$\times$10$^{-3}$ (2.57$\times$10$^{-3}$) eV$^{2}$/c$^{4}$ 
and 
0.433 (0.456) $< \sin^{2}\theta_{23} <$ 0.594 (0.572)
for the normal mass hierarchy,
and
2.37$\times$10$^{-3}$ (2.42$\times$10$^{-3}$) eV$^{2}$/c$^{4}$ 
$< |\Delta m^{2}_{32}| < 
$ 2.70$\times$10$^{-3}$ (2.64$\times$10$^{-3}$) eV$^{2}$/c$^{4}$ 
and 
0.431 (0.454) $< \sin^{2}\theta_{23} <$ 0.591 (0.569)
for the inverted mass hierarchy.

\end{abstract}

%
% Keywords
%

%\begin{keyword}
%\end{keyword}

\end{frontmatter}

\clearpage

\tableofcontents

\clearpage

\listoffigures
\listoftables


